\section{Роботи і люди}

% ===== Дорожня карта + Ґант =====
% Місяці робіт — з gantt.pdf (Roadmap.md), місяць 1 = листопад 2026 (RACI.md).
% TODO: у gantt.pdf роботи MVP (фази 2–4) стоять у місяцях 3–10, а в Project Overview
% MVP — 4–20 листопада, бета — грудень–травень. Коли команда узгодить Ґант
% з Overview, оновити місяці нижче.
\begin{frame}{Дорожня карта: 25 робіт у 5 фазах}
\vspace{-0.1cm}
\centering
\begin{tikzpicture}[x=1.183cm,y=1cm]
  \foreach \m [count=\i] in {лис,гру,січ,лют,бер,кві,тра,чер,лип,сер,вер,жов}{
    \node[font=\scriptsize\semibold,text=ink] at (\i-0.5,0.18) {\i};
    \node[font=\tiny,text=mute] at (\i-0.5,-0.15) {\m};
  }
  \foreach \i in {0,...,12}{\draw[line,line width=0.4pt] (\i,0.38) -- (\i,-4.75);}
  \gphase{1}{2}{Дослідження та підготовка}{1}
  \gphase{3}{5}{Обробка даних (ядро)}{2}
  \gphase{6}{7}{Картографія та аналітика}{3}
  \gphase{8}{10}{Рушій порад та інтерфейс}{4}
  \gphase{11}{12}{Тестування та реліз}{5}
  % Фаза 1
  \gwork{1}{2}{1}{1}{1}{1. Юридичний аудит карт}
  \gwork{1}{2}{1}{2}{2}{2. Проєктування архітектури}
  \gwork{1}{2}{1}{2}{3}{3. UI/UX дизайн}
  \gwork{1}{2}{2}{2}{4}{4. Налаштування CI/CD}
  \gwork{1}{2}{1}{1}{5}{5. Технічна валідація GPMF}
  % Фаза 2
  \gwork{3}{5}{3}{4}{1}{6. Модуль парсингу відео}
  \gwork{3}{5}{3}{4}{2}{7. Алгоритм поділу на кола}
  \gwork{3}{5}{4}{5}{3}{8. Фільтрація невалідних кіл}
  \gwork{3}{5}{4}{5}{4}{9. Корекція зсуву GPS}
  \gwork{3}{5}{3}{5}{5}{10. Серверна інфраструктура}
  % Фаза 3
  \gwork{6}{7}{6}{6}{1}{11. Оцифрування траси MVP}
  \gwork{6}{7}{6}{6}{2}{12. Рендеринг траєкторії}
  \gwork{6}{7}{6}{7}{3}{13. Ідентифікація поворотів}
  \gwork{6}{7}{6}{7}{4}{14. Механізм порівняння кіл}
  \gwork{6}{7}{7}{7}{5}{15. Сортування втрат часу}
  % Фаза 4
  \gwork{8}{10}{8}{9}{1}{16. Рушій текстових порад}
  \gwork{8}{10}{8}{9}{2}{17. Синхронізація з відео}
  \gwork{8}{10}{9}{10}{3}{18. Генератор плану тренувань}
  \gwork{8}{10}{8}{10}{4}{19. Верстка дашборда}
  \gwork{8}{10}{10}{10}{5}{20. Система оцінювання порад}
  % Фаза 5
  \gwork{11}{12}{11}{11}{1}{21. Точнісне тестування}
  \gwork{11}{12}{11}{12}{2}{22. Закрите бета-тестування}
  \gwork{11}{12}{11}{11}{3}{23. Налаштування продуктової аналітики}
  \gwork{11}{12}{12}{12}{4}{24. Валідація бізнес-моделі}
  \gwork{11}{12}{12}{12}{5}{25. Підготовка до масштабування}
\end{tikzpicture}

\vspace{0.05cm}
\raggedright{\tiny\color{mute}Місяць 1 — листопад 2026. Повна діаграма Ґанта — gantt.pdf, опис кожної роботи — Roadmap.md.\par}
\note{\speaker{П2 · Дарина}
Дорожня карта — це 25 робіт у п'яти фазах, які йдуть одна за одною протягом року.

Перші два місяці — дослідження й підготовка: юридичний аудит карт, архітектура, дизайн, CI/CD і перевірка, що бібліотека gopro-telemetry стабільно дістає GPS з відео. Далі три місяці на ядро обробки даних: парсинг відео, поділ на кола, фільтрація і сервер. Потім картографія й аналітика, рушій порад з інтерфейсом, а останні два місяці — тестування точності, закрита бета й аналітика.

Кожна червона смужка — робота на діаграмі Ґанта. Повна діаграма — у файлі gantt.pdf.}
\end{frame}

% ===== RACI =====
% Дані — RACI.md. Рядки — учасниці, стовпці — роботи 1–25 з дорожньої карти.
\begin{frame}{RACI: хто за що відповідає}
\vspace{0.05cm}
\centering
\begin{tikzpicture}[x=1cm,y=1cm]
  % фази і номери робіт
  \foreach \s/\e/\n in {1/5/1,6/10/2,11/15/3,16/20/4,21/25/5}{
    \draw[ink,line width=0.6pt] ({(\s-1)*0.42+0.03},0.78) -- ({\e*0.42-0.03},0.78);
    \node[font=\tiny\semibold,text=ink] at ({(\s+\e-1)*0.21},0.92) {Ф\n};
  }
  \foreach \c in {1,...,25}{\node[font=\tiny\ttfamily,text=mute] at ({(\c-1)*0.42+0.21},0.5) {\c};}
  \node[font=\tiny\semibold,text=ink] at (10.95,0.5) {R};
  \node[font=\tiny\semibold,text=kerb] at (11.5,0.5) {A};
  \racirow{1}{Демянчук Д.}{бек/тест}{I,AR,I,R,AR,C,C,R,C,AR,I,I,AR,C,I,R,C,C,R,AR,C,C,C,I,R}
  \racirow{2}{Наливайко К.}{аналітик/бек}{C,C,C,I,I,I,AR,C,AR,I,C,C,R,AR,C,AR,I,C,I,I,R,C,AR,C,R}
  \racirow{3}{Проніна М.}{фронт}{I,C,AR,I,C,AR,I,-,-,C,R,AR,I,-,-,I,AR,I,AR,C,I,C,R,C,C}
  \racirow{4}{Афанасенко Я.}{бек/тест}{I,C,I,AR,C,I,R,AR,C,C,I,I,C,R,AR,I,R,AR,I,R,AR,R,I,I,C}
  \racirow{5}{Уточкіна Я.}{лід}{AR,I,R,I,I,I,I,I,I,I,AR,I,I,-,-,R,I,R,C,C,I,AR,I,AR,A}
  % підсумки
  \foreach \r/\R in {1/10,2/8,3/7,4/10,5/7}{
    \node[font=\scriptsize\heavy,text=ink] at (10.95,{-(\r-1)*0.62}) {\R};
    \node[font=\scriptsize\heavy,text=kerb] at (11.5,{-(\r-1)*0.62}) {5};
  }
  \draw[line] (10.65,0.35) -- (10.65,-2.8);
\end{tikzpicture}

\vspace{0.25cm}
{\tiny\color{body}
\tikz[baseline=-0.6ex]\fill[ink] (0,-0.1) rectangle (0.2,0.12);\;\textbf{R} — виконує\qquad
\tikz[baseline=-0.6ex]\fill[kerb] (0,-0.1) rectangle (0.2,0.12);\;\textbf{A/R} — виконує і відповідає\qquad
\tikz[baseline=-0.6ex]\draw[kerb,line width=0.8pt] (0,-0.1) rectangle (0.2,0.12);\;\textbf{A} — відповідає, рівно одна на роботу\qquad
\tikz[baseline=-0.6ex]\fill[line] (0,-0.1) rectangle (0.2,0.12);\;\textbf{C} — радяться\qquad
{\color{later}\textbf{I}} — інформують}
\note{\speaker{П4 · Ярина}
RACI-матриця: рядки — це ми п'ятеро, стовпці — 25 робіт з дорожньої карти, з тими самими номерами.

Правила такі: у кожній роботі рівно одна людина відповідає за результат — це A, і хоча б одна виконує — це R. Червоні клітинки — коли одна людина і виконує, і відповідає.

Праворуч — підсумки. Кожна з нас відповідає за п'ять робіт. Виконавицею найчастіше стоять Дарина і Ярина — по 10 робіт, Катерина — 8, Марина і Яна — по 7.}
\end{frame}

% ===== Навантаження =====
% Дані — RACI.md, розділ «Навантаження по місяцях».
\begin{frame}{Навантаження по місяцях}
\vspace{0.1cm}
\begin{columns}[T,onlytextwidth]
\column{0.68\textwidth}
\begin{tikzpicture}
  \foreach \m [count=\i] in {лис,гру,січ,лют,бер,кві,тра,чер,лип,сер,вер,жов}{
    \node[font=\tiny\ttfamily,text=mute] at ({(\i-0.5)*0.62},0.62) {\i};
    \node[font=\tiny,text=mute] at ({(\i-0.5)*0.62},0.4) {\m};}
  \loadrow{1}{Демянчук Д.}{2,2,1,2,2,1,1,2,2,2,0,1}
  \loadrow{2}{Наливайко К.}{0,0,1,2,1,2,2,1,1,0,2,1}
  \loadrow{3}{Проніна М.}{1,1,1,1,0,2,0,2,2,1,1,0}
  \loadrow{4}{Афанасенко Я.}{0,1,1,2,1,1,2,1,2,2,2,1}
  \loadrow{5}{Уточкіна Я.}{2,1,0,0,0,1,0,1,2,1,1,2}
\end{tikzpicture}

\vspace{0.2cm}
{\tiny\color{mute}Кількість одночасних робіт, де людина R. Місяць 1 — листопад 2026.\par}
\column{0.29\textwidth}
  \lbl{Перевірка матриці}\par\vskip0.2cm
  {\scriptsize
  \ok У кожній роботі рівно одна A\par\vskip0.12cm
  \ok У кожній роботі хоча б одна R\par\vskip0.12cm
  \ok R пораховано на кожну\par\vskip0.12cm
  \ok Ніхто не має більше 2 R одночасно}
\end{columns}
\note{\speaker{П4 · Ярина}
Щоб ніхто не перевантажився, ми розклали RACI по місяцях за діаграмою Ґанта. У клітинці — скільки робіт людина виконує одночасно в цьому місяці.

Максимум — дві роботи одночасно, більше ні в кого немає. Ми також перевірили, що в кожній роботі рівно одна A і хоча б одна R, і порахували R на кожну учасницю.}
\end{frame}
